% GENERATED FILE. Edit canonical lessons, then run npm run generate:books.
% canonical-source-hash: fnv1a64:68b57a1a393631f0
% canonical-lessons: ZH-C65-men, ZH-W65-men-door, ZH-C65-qingwen, ZH-W65-wen-ask, ZH-C65-buhaoyisi, ZH-C65-kai, ZH-C65-guan

\chapter{Excuse Me}
\label{ch:zh-excuse-me}
\hlchaptermodality{65}

\begin{chapteropening}
\textbf{By the end of this chapter:} \emph{I can say excuse me, ask politely, and ask somebody to open or shut something.}

Complete the last lesson of chapter 65: I can say excuse me, ask politely, and ask somebody to open or shut something.
\end{chapteropening}

\section[\texorpdfstring{\zh{门} (mén)}{门 (mén)}]{\zh{门} (mén) — door; gate}

\label{lesson:ZH-C65-men}



\begin{quote}
Before the new one: say the Mandarin for to borrow, then the Mandarin for to give back.
\end{quote}

\subsection*{You'll want to know: \zh{门}}
\textbf{\zh{门}} — \emph{mén} — \textquotedblleft{}door; gate\textquotedblright{}.

\begin{quote}
\textbf{\zh{这是我家的门}} \emph{zhè shì wǒ jiā de mén} — This is the door to my home.
\end{quote}

The new character, in pieces: \zh{门} — a dot, a post down the left, and a frame that turns down the right and hooks: the right-hand piece of \zh{们}. The next lesson puts it in your hand.

\begin{cousinweb}
\textbf{\zh{门}} (\emph{mén}) is inherited from Middle Chinese \emph{mwon} (Guangyun, in Baxter's transcription); the dictionary says its old character is made of two door leaves.
\end{cousinweb}

\begin{grammarlens}[title={What you've built}]
One more word for asking politely.
\end{grammarlens}

\subsection*{Your turn}
\begin{itemize}
\raggedright
  \item \emph{Say it:} \emph{mén}
  \item \emph{Say it:} \emph{mén}, once more
  \item \emph{Say it:} \emph{zhè shì wǒ jiā de mén} — This is the door to my home.
  \item \emph{Recall:} say the Mandarin for to borrow, then the Mandarin for to give back, then say \emph{mén} again
\end{itemize}

\begin{tcolorbox}[breakable,colback=teal!4,colframe=teal!35!black,title={Before you move on}]
What does \zh{门} mean? (\textquotedblleft{}Door; gate\textquotedblright{}.) Say it once more.
\end{tcolorbox}
\section[\texorpdfstring{\zh{门} (mén)}{门 (mén)}]{\zh{门} — a dot, a post, a frame}

\label{lesson:ZH-W65-men-door}



\begin{quote}
The word first, then the character inside it.

\begin{itemize}
\raggedright
  \item \emph{Recall:} say \emph{mén} — \textbf{R1}, one lesson back, before it decays
\end{itemize}

\begin{quote}
\zh{们} \zh{口}
\end{quote}

\emph{Read it:} each one aloud, with nothing said first — \emph{men}, \emph{kǒu}
\end{quote}

\subsection*{Script — a dot, a post, a frame}
\hlblockfigure{\detokenize{figures/ZH-W65-men-door-filmstrip.pdf}}{How \zh{门} is written, stroke by stroke}

\begin{quote}
\zh{门}
\end{quote}

\textbf{One piece, three strokes}, the same three as on the right of \zh{们}.

\textbf{A dot} at the top left.

\textbf{A vertical} down the left side, the door post.

\textbf{A horizontal that turns down the right side} and hooks to the left at the foot.

\subsection*{Your turn}
\begin{itemize}
\raggedright
  \item \emph{Write it:} \zh{门}, all three, with the model in front of you
  \item \emph{Cover:} write it again from memory, and count the pieces, not the strokes
\end{itemize}

\begin{tcolorbox}[breakable,colback=teal!4,colframe=teal!35!black,title={Before you move on}]
How many strokes in \zh{门}? (\textbf{Three}.) Which character you know has it on the right? (\textbf{\zh{们}}.) How does the last stroke end? (\textbf{It hooks to the left}.)

Source: \href{https://raw.githubusercontent.com/chanind/hanzi-writer-data/68d10a4b21150cae5e1ebbd223eed289cf32d90c/data/\%E9\%97\%A8.json}{Hanzi Writer Data \zh{门}, pinned stroke paths and medians}.
\end{tcolorbox}
\section[\texorpdfstring{\zh{请问} (qǐngwèn)}{请问 (qǐngwèn)}]{\zh{请问} (qǐngwèn) — excuse me, may I ask}

\label{lesson:ZH-C65-qingwen}



\begin{quote}
Before the new one: say the Mandarin for to give back, then the Mandarin for door.
\end{quote}

\subsection*{You'll want to know: \zh{请问}}
\textbf{\zh{请问}} — \emph{qǐngwèn} — \textquotedblleft{}excuse me, may I ask\textquotedblright{}.

\begin{quote}
\textbf{\zh{请问这是你的手机吗}} \emph{qǐngwèn zhè shì nǐ de shǒujī ma} — Excuse me, is this your phone?
\end{quote}

The new character, in pieces: \zh{问} — \zh{门} with \zh{口} inside it. The next lesson puts it in your hand.

\begin{cousinweb}
\textbf{\zh{请问}} (\emph{qǐngwèn}) is built from \zh{请} 'please' and \zh{问} 'to ask', both inherited from Middle Chinese, \emph{tshjengX} and \emph{mjunH} (Guangyun, in Baxter's transcription); the dictionary already glosses the first as 'to beg, to ask'.
\end{cousinweb}

\begin{grammarlens}[title={What you've built}]
One more word for asking politely.
\end{grammarlens}

\subsection*{Your turn}
\begin{itemize}
\raggedright
  \item \emph{Say it:} \emph{qǐngwèn}
  \item \emph{Say it:} \emph{qǐngwèn}, once more
  \item \emph{Say it:} \emph{qǐngwèn zhè shì nǐ de shǒujī ma} — Excuse me, is this your phone?
  \item \emph{Recall:} say the Mandarin for to give back, then the Mandarin for door, then say \emph{qǐngwèn} again
\end{itemize}

\begin{tcolorbox}[breakable,colback=teal!4,colframe=teal!35!black,title={Before you move on}]
What does \zh{请问} mean? (\textquotedblleft{}Excuse me, may I ask\textquotedblright{}.) Say it once more.
\end{tcolorbox}
\section[\texorpdfstring{\zh{问} (wèn)}{问 (wèn)}]{\zh{问} — \zh{门} with \zh{口} inside}

\label{lesson:ZH-W65-wen-ask}



\begin{quote}
The word first, then the character inside it.

\begin{itemize}
\raggedright
  \item \emph{Recall:} say \emph{qǐngwèn} — \textbf{R1}, one lesson back, before it decays
\end{itemize}

\begin{quote}
\zh{门} \zh{口}
\end{quote}

\emph{Read it:} each one aloud, with nothing said first — \emph{mén}, \emph{kǒu}
\end{quote}

\subsection*{Script — \zh{门} with \zh{口} inside}
\hlblockfigure{\detokenize{figures/ZH-W65-wen-ask-filmstrip.pdf}}{How \zh{问} is written, stroke by stroke}

\begin{quote}
\zh{问}
\end{quote}

\textbf{Two pieces: the door first, then the mouth inside it.}

\textbf{\zh{门}, three strokes}: the dot, the left post, the frame that turns down and hooks.

\textbf{\zh{口}, three strokes}, small and inside: the left side, the top that turns down, the bottom that closes it.

3 + 3 = six, and both pieces are already yours.

\subsection*{Your turn}
\begin{itemize}
\raggedright
  \item \emph{Write it:} \zh{问}, all six, with the model in front of you
  \item \emph{Cover:} write it again from memory, and count the pieces, not the strokes
\end{itemize}

\begin{tcolorbox}[breakable,colback=teal!4,colframe=teal!35!black,title={Before you move on}]
How many pieces in \zh{问}? (\textbf{Two} — \zh{门}, then \zh{口}.) Which goes first? (\textbf{\zh{门}, the frame}.) Where does \zh{口} sit? (\textbf{Inside the door}.)

Source: \href{https://raw.githubusercontent.com/chanind/hanzi-writer-data/68d10a4b21150cae5e1ebbd223eed289cf32d90c/data/\%E9\%97\%AE.json}{Hanzi Writer Data \zh{问}, pinned stroke paths and medians}.
\end{tcolorbox}
\section[bù hǎoyìsi]{bù hǎoyìsi — excuse me; sorry to trouble you}

\label{lesson:ZH-C65-buhaoyisi}



\begin{quote}
Before the new one: say the Mandarin for door, then the Mandarin for may I ask.
\end{quote}

\subsection*{You'll want to know: bù hǎoyìsi}
\emph{bù hǎoyìsi} — \textquotedblleft{}excuse me; sorry to trouble you\textquotedblright{}.

\begin{quote}
\emph{bù hǎoyìsi wǒ xiǎng wèn nǐ} — Excuse me, could I ask you something?
\end{quote}

Its characters come later; for now it is a word for the ear.

bù keeps its own falling tone here, because hǎo after it is a third tone, not a fourth.

\begin{cousinweb}
\emph{bù hǎoyìsi} is built from bù 'not', hǎo 'good' and yìsi 'meaning', all inherited from Middle Chinese, \emph{pjut}, \emph{xawX}, \emph{'iH} and \emph{si} (Guangyun, in Baxter's transcription); the dictionary glosses the last two as 'intention' and 'to think of'. 'Not a good feeling' is the literal sense; how it became an apology is not settled.
\end{cousinweb}

\begin{grammarlens}[title={What you've built}]
One more word for asking politely.
\end{grammarlens}

\subsection*{Your turn}
\begin{itemize}
\raggedright
  \item \emph{Say it:} \emph{bù hǎoyìsi}
  \item \emph{Say it:} \emph{bù hǎoyìsi}, once more
  \item \emph{Say it:} \emph{bù hǎoyìsi wǒ xiǎng wèn nǐ} — Excuse me, could I ask you something?
  \item \emph{Recall:} say the Mandarin for door, then the Mandarin for may I ask, then say \emph{bù hǎoyìsi} again
\end{itemize}

\begin{tcolorbox}[breakable,colback=teal!4,colframe=teal!35!black,title={Before you move on}]
What does bù hǎoyìsi mean? (\textquotedblleft{}Excuse me; sorry to trouble you\textquotedblright{}.) Say it once more.
\end{tcolorbox}
\section[kāi]{kāi — to open; to switch on}

\label{lesson:ZH-C65-kai}



\begin{quote}
Before the new one: say the Mandarin for may I ask, then the Mandarin for excuse me.
\end{quote}

\subsection*{You'll want to know: kāi}
\emph{kāi} — \textquotedblleft{}to open; to switch on\textquotedblright{}.

\begin{quote}
\emph{shāngdiàn jiǔ diǎn kāi} — The shop opens at nine.
\end{quote}

Its characters come later; for now it is a word for the ear.

\begin{cousinweb}
\emph{kāi} is inherited from Middle Chinese \emph{khoj} (Guangyun, in Baxter's transcription), glossed 'to open, to undo'.
\end{cousinweb}

\begin{grammarlens}[title={What you've built}]
One more word for asking politely.
\end{grammarlens}

\subsection*{Your turn}
\begin{itemize}
\raggedright
  \item \emph{Say it:} \emph{kāi}
  \item \emph{Say it:} \emph{kāi}, once more
  \item \emph{Say it:} \emph{shāngdiàn jiǔ diǎn kāi} — The shop opens at nine.
  \item \emph{Recall:} say the Mandarin for may I ask, then the Mandarin for excuse me, then say \emph{kāi} again
\end{itemize}

\begin{tcolorbox}[breakable,colback=teal!4,colframe=teal!35!black,title={Before you move on}]
What does kāi mean? (\textquotedblleft{}To open; to switch on\textquotedblright{}.) Say it once more.
\end{tcolorbox}
\section[guān]{guān — to shut; to close; to switch off}

\label{lesson:ZH-C65-guan}



\begin{quote}
Before the new one: say the Mandarin for excuse me, then the Mandarin for to open.
\end{quote}

\subsection*{You'll want to know: guān}
\emph{guān} — \textquotedblleft{}to shut; to close; to switch off\textquotedblright{}.

\begin{quote}
\emph{qǐng guān mén} — Please shut the door.
\end{quote}

Its characters come later; for now it is a word for the ear.

\begin{cousinweb}
\emph{guān} is inherited from Middle Chinese \emph{kwaen} (Guangyun, in Baxter's transcription), which the dictionary explains, through an older book, as a wooden bar held across a door. Shutting is probably what the bar does, but the step is not given there.
\end{cousinweb}

\begin{grammarlens}[title={What you've built}]
One more word for asking politely.
\end{grammarlens}

\subsection*{Your turn}
\begin{itemize}
\raggedright
  \item \emph{Say it:} \emph{guān}
  \item \emph{Say it:} \emph{guān}, once more
  \item \emph{Say it:} \emph{qǐng guān mén} — Please shut the door.
  \item \emph{Recall:} say the Mandarin for excuse me, then the Mandarin for to open, then say \emph{guān} again
\end{itemize}

\begin{tcolorbox}[breakable,colback=teal!4,colframe=teal!35!black,title={Before you move on}]
What does guān mean? (\textquotedblleft{}To shut; to close; to switch off\textquotedblright{}.) Say it once more.
\end{tcolorbox}
